\BourbakiExercises{习题}

\begin{enumerate}
\def\labelenumi{\arabic{enumi})}
\tightlist
\item
  设 A 与 B 是 K-代数，M 是一个 A-模；把 \(M \otimes B\) 视为代数 \(C = A \otimes B\) 上的模。把 M 等同为 \(M \otimes B\) 的一个 A-子模。
\end{enumerate}

\begin{enumerate}
\def\labelenumi{\alph{enumi})}
\item
  证明包含关系 \(M \cap \mathfrak{R}_{\mathrm{C}}(\mathrm{M} \otimes \mathrm{B}) \subset \mathfrak{R}_{\mathrm{A}}(\mathrm{M})\)（考虑一个从 M 到某个单 A-模的 A-线性同态）。
\item
  证明在下列每种情形中等号成立：
\end{enumerate}

\begin{enumerate}
\def\labelenumi{(\roman{enumi})}
\item
  A 在 \(\mathbf{K}\) 上是有限维的。
\item
  A-模 M 是有限生成的，且 A 是一个由在 K 上有限维的子代数组成的右向定向族的并。
\item
  \(M = A_{s}\)，且 A 的根是幂零的。
\end{enumerate}

（在情形 (i) 中，可证明每个从 \(\mathrm{M} \otimes \mathrm{B}\) 到单 C-模 S 的同态都在 \(\Re_{\mathrm{A}}(\mathrm{M})\) 上为零，办法是注意到 \(\Re_{\mathrm{C}}(\mathrm{S})\) 为零。）

\begin{enumerate}
\def\labelenumi{\arabic{enumi})}
\setcounter{enumi}{1}
\item
  设 \(\mathrm{A}\) 是一个为整环的 K-代数，\(\mathrm{L}\) 是它的分式域。证明 \(\mathrm{L}\) 是一个单 \(\mathrm{A} \otimes \mathrm{L}\) -模，但它作为 A-模的根等于 \(\mathrm{L}\)。
\item
  给出（非零）K-代数 \(\mathrm{A}_1\) 与 \(\mathrm{A}_2\) 的一个例子，使得 \(\Re (\mathrm{A}_1)\neq 0\) 而 \(\Re (\mathrm{A}_1\otimes \mathrm{A}_2) = 0\)（参见 VIII, 第 167 页，习题 11）。
\item
  给出一个交换域 K 与 K 的两个扩张 E 和 F 的例子，使得环 \(A = E \otimes_{K} F\) 是半单的但不是域（相应地，不是半单的）。推出 A-模 \(A_{s} = E_{s} \otimes_{K} F_{s}\) 是半单的但不是单的（相应地，不是半单的）。
\end{enumerate}

¶ 5) 设 A 是一个拟单 K-代数（VIII, 第 120 页，注 2）。

\begin{enumerate}
\def\labelenumi{\alph{enumi})}
\item
  证明 (A, A)-双模 \(_{s}A_{d}\) 是单的；推出 A 的中心 Z 是一个体，同构于该双模的交换子。
\item
  设 B 是一个 K-代数。证明 \(A \otimes B\) 的每个双边理想都具有形式 \(A \otimes_{Z} b\)，其中 b 是 B 的一个双边理想（化归到情形 Z = K，并利用 VIII, 第 63 页的推论 2）。
\item
  推出若 A 与 B 是拟单 K-代数且其中之一的中心为 K，则代数 \(\mathrm{A} \otimes \mathrm{B}\) 是拟单的。
\end{enumerate}

\(d\) ) 设 B 是一个 K-代数，A 是 B 的一个中心 Z 包含 K 的拟单子代数。证明 A 与它在 B 中的交换子 \(\mathrm{A}'\) 在 Z 上线性分离（利用 \(b\)）。

\begin{enumerate}
\def\labelenumi{\alph{enumi})}
\setcounter{enumi}{4}
\item
  设 A 是一个拟单 K-代数，Z 是它的中心。设 \(\varphi : \mathrm{A} \otimes \mathrm{A}^{\circ} \to \operatorname{End}_{\mathrm{Z}}(\mathrm{A})\) 是把 \(a \otimes b\) 映为自同态 \(x \mapsto axb\) 的同态。证明 \(\varphi\) 是单射且它的像 H 是一个拟单代数（利用 \(c\) 与 \(d\)）。此外证明 H 在 \(\operatorname{End}_{\mathrm{Z}}(\mathrm{A})\) 中稠密（VIII, 第 129 页，习题 3）；我们有 \(\mathrm{H} = \operatorname{End}_{\mathrm{Z}}(\mathrm{A})\) 当且仅当 A 在 Z 上是有限维的（参见 VIII, 第 128 页，习题 2）。
\item
  设 A 是一个以 K 为中心的半单代数。由 e) 推出，K-代数 \(A \otimes A^{o}\) 是半单的当且仅当 A 在 K 上秩有限（参见 VIII, 第 238 页，定理 2）。
\end{enumerate}

\begin{enumerate}
\def\labelenumi{\arabic{enumi})}
\setcounter{enumi}{5}
\item
  设 A 是一个具有非零幂零根的 K-代数。证明同态 \(\varphi: A \otimes A^{\circ} \to \operatorname{End}_{K}(A)\)（习题 5）不是单射。
\item
  设 \(A_{1}\) 与 \(A_{2}\) 是半单 K-代数，\(Z_{1}\) 与 \(Z_{2}\) 是它们的中心。假设 \(A_{1}\) 在 K 上次数有限，且环 \(Z_{1} \otimes Z_{2}\) 是约化的。设 f 是一个从 \(A_{2}\) 到 \(A_{1}\) 的 K-代数同态。证明 \(f(\mathrm{A}_{2})\) 在 \(A_{1}\) 中的交换子是一个半单 K-代数（把 \(A_{1}\) 视为 \((\mathrm{A}_{2} \otimes \mathrm{A}_{1}^{\circ})\) -模，并利用 VIII, 第 221 页的命题 7）。
\item
  设 K 是一个交换域，L 是 K 的一个次数 \textgreater{} 1 的有限根式扩张。设 V 是 L 上的一个无限维向量空间，A 是 \(\operatorname{End}_{\mathrm{L}}(\mathrm{V})\) 的由有限秩自同态与恒等映射生成的 K-子代数。证明 L-代数 \(\mathrm{A}_{(\mathrm{L})}\) 具有非零根，尽管其中心等于 L（利用习题 7）。
\item
  设 A 与 B 是 K-代数，M 是一个自由 A-模，N 是一个 B-模。
\end{enumerate}

\begin{enumerate}
\def\labelenumi{\alph{enumi})}
\item
  证明典范同态 \(\vartheta\)（从 \(\mathrm{End}_{\mathrm{A}}(\mathrm{M})\otimes \mathrm{End}_{\mathrm{B}}(\mathrm{N})\) 到 \(\mathrm{End}_{\mathrm{A}\otimes \mathrm{B}}(\mathrm{M}\otimes \mathrm{N})\)）的像是环 \(\mathrm{End}_{\mathrm{A}\otimes \mathrm{B}}(\mathrm{M}\otimes \mathrm{N})\) 的一个稠密子环。
\item
  若 B-模 N 的位似环在其双交换子中稠密，则 \((A \otimes B)\) -模 \(M \otimes N\) 的位似环亦然。
\item
  假设存在 N 的一个元素 y 与 N 的自同态的一个序列 \((v_{n})\)，使得诸元素 \(v_{n}(y)\) 生成 K 上的一个无限维向量空间。证明映射 \(\vartheta\) 不是满射。
\end{enumerate}

\begin{enumerate}
\def\labelenumi{\arabic{enumi})}
\setcounter{enumi}{9}
\tightlist
\item
  保留习题 9 的假设，并此外假设 B-模 N 是单的；用 D 表示它的交换子。
\end{enumerate}

\begin{enumerate}
\def\labelenumi{\alph{enumi})}
\item
  映射 \(\vartheta\) 是满射的当且仅当 M 在 A 上具有有限基，或者 D 在 K 上是有限维的（利用习题 9, c)）。
\item
  证明 \((\mathrm{A} \otimes \mathrm{B})\) -模 \(M \otimes N\) 是半单的当且仅当环 \(A \otimes D^{o}\) 是半单的（化归到情形 \(M = A_{s}\)；注意到 \(M \otimes N\) 是一个自由 \(A \otimes D^{o}\) -模，并利用 VIII, 第 85 页的命题 5 以及习题 9, a)）。
\end{enumerate}

¶ 11) 设 A 与 B 是 K-代数；假设环 A 与 B 都是本原的，且它们各自的底座 s 与 t（VIII, 第 130 页，习题 8）不化为 0。设 D（相应地，E）是 A（相应地，B）的一个极小左理想 a（相应地，b）的交换子；假设代数 D ⊗ E 是单的。证明 A ⊗ B 是一个以 s ⊗ t 为底座的本原环（注意到 (A ⊗ B)-模 a ⊗ b 是有限长度的同型模；利用 VIII, 第 131 页的习题 9, a)）。

\begin{enumerate}
\def\labelenumi{\arabic{enumi})}
\setcounter{enumi}{11}
\tightlist
\item
  设 \(A_{1}\) 是 K-代数 \(\mathrm{K}(\mathrm{X})[\mathrm{U}]_{1,d/d\mathrm{X}}\)（VIII, 第 11 页），\(A_{2}\) 是 K-代数 \(K[[T]]\)，A 是 K-代数 \(A_{1} \otimes_{K} A_{2}\)。把 \(A_{1}\) 与 \(A_{2}\) 等同为 A 的子代数。
\end{enumerate}

\begin{enumerate}
\def\labelenumi{\alph{enumi})}
\item
  证明 \(A_{2}\) 是 A 的中心，A 不含任何零因子，且 A 的可逆元恰是 \(A_{2}\) 的那些元素。
\item
  证明 A 的双边理想恰是诸理想 \(AT^{n}\)（\(n \geqslant 0\)）。
\item
  证明环 A 是本原的，尽管其中心具有非零根（注意到 A 的一个包含 1 - TU 的极大左理想不含 0 以外的任何双边理想）。
\end{enumerate}

\begin{enumerate}
\def\labelenumi{\arabic{enumi})}
\setcounter{enumi}{12}
\tightlist
\item
  设 D 与 E 是中心包含 K 的体，V（相应地，W）是 D（相应地，E）上的一个向量空间。证明 \((\mathrm{D} \otimes_{\mathrm{K}} \mathrm{E})\) -模 \(V \otimes_{K} W\) 是自由的。
\end{enumerate}

¶ 14) 设 V 是一个体 D 上的向量空间。用 \(\Omega\) 表示加群 V 的自同态环，并把 D 等同为 \(\Omega\) 的一个子环。令 \(A = \text{End}_{\text{D}}(\text{V})\)。

\begin{enumerate}
\def\labelenumi{\alph{enumi})}
\item
  证明 A 与 D 在它们的公共中心 Z 上线性分离（参见习题 5, d)）。
\item
  \(D \otimes_{Z} A\) 在 \(\Omega\) 中的像等于 \(\operatorname{End}_{\mathbb{Z}}(\mathrm{V})\) 当且仅当 D 在 Z 上秩有限。（为看出该条件是必要的，首先由习题 5, e) 推出 V = D 的情形，然后固定 V 的一个元素 \(x \neq 0\)，并考虑 \(\operatorname{End}_{\mathbb{Z}}(\mathrm{V})\) 的子代数 B，它由满足 \(u(\mathrm{D}x) \subset \mathrm{D}x\) 的 V 的 Z-自同态 u 组成。注意到等式 \(\operatorname{End}_{\mathbb{Z}}(\mathrm{V}) = \mathrm{DA}\) 蕴含 \(\mathrm{B} = \mathrm{D}(\mathrm{B} \cap \mathrm{A})\)。）
\end{enumerate}

\begin{enumerate}
\def\labelenumi{\arabic{enumi})}
\setcounter{enumi}{14}
\tightlist
\item
  \begin{enumerate}
  \def\labelenumii{\alph{enumii})}
  \tightlist
  \item
    设 \(\mathrm{K}\) 是一个交换域，A 是一个 K-代数，\(\mathrm{B} = \mathrm{K}[\mathrm{X}_{\iota}]_{\iota \in \mathrm{I}}\) 是 \(\mathrm{K}\) 上的一个多项式代数。证明若 \(\mathrm{A}\) 的每个非零元素都是正则的，则代数 \(\mathrm{A} \otimes_{\mathrm{K}} \mathrm{B}\) 具有同样的性质。
  \end{enumerate}
\end{enumerate}

\begin{enumerate}
\def\labelenumi{\alph{enumi})}
\setcounter{enumi}{1}
\tightlist
\item
  推出若 D 是一个在 K 上秩有限的体，而 E 是 K 的一个纯超越扩张，则环 \(D \otimes_{K} E\) 是一个体。
\end{enumerate}

¶ 16) 设 D 是一个体，V 是 D 上的一个无限维右向量空间，A 是环 \(\mathrm{End}_{\mathrm{D}}(\mathrm{V})\) 的一个底座 \(\mathfrak{s}\) 非零的稠密子环（VIII, 第 130 页，习题 8 与 9）。设 K 是 D 的中心的一个子域，B 是一个 K-代数。证明若 K-代数 \(\mathrm{D}^{\circ} \otimes \mathrm{B}\) 的根非零，则 \(\mathrm{A} \otimes \mathrm{B}\) 的根亦非零。（设 \(\mathfrak{a}\) 为元素 \(\sum u_{i} \otimes b_{i}\)（属于 \(\mathfrak{s} \otimes \mathrm{B}\)）中满足 \(\sum \langle u_{i}(x), x^{*} \rangle \otimes b_{i} \in \Re(\mathrm{D}^{\circ} \otimes \mathrm{B})\)（对所有 \(x \in \mathrm{V}\) 与 \(x^{*} \in \mathrm{V}^{*}\)）者所成之集。证明 \(\mathfrak{a}\) 是 \(\mathrm{A} \otimes \mathrm{B}\) 的一个双边理想，并利用 VIII, 第 167 页的习题 10, e) 与 VIII, 第 131 页的习题 9，证明它非零且包含在 \(\mathrm{A} \otimes \mathrm{B}\) 的根中。）

\begin{enumerate}
\def\labelenumi{\arabic{enumi})}
\setcounter{enumi}{16}
\tightlist
\item
  设 K 是一个交换域，D 与 E 是中心包含 K 的体。假设 \([E:K]=m\) 有限，且体 D 与 E 之一的中心等于 K。则 K-代数 \(D\otimes_{K}E\) 同构于一个体 C 上的有限维向量空间 h 的自同态代数（VIII, 第 222 页，推论 2）。
\end{enumerate}

\begin{enumerate}
\def\labelenumi{\alph{enumi})}
\item
  证明 \(h\) 整除 \(m\)（参见 VIII, 第 205 页，公式 (32)）。
\item
  设 W 是 D 上的一个向量空间；则 E 同构于 \(\operatorname{End}_{\mathrm{D}}(\mathrm{W})\) 的一个 K-子代数当且仅当 W 的维数是无限的或是 m/h 的倍数（注意到此时 W 是一个 \((\mathrm{D} \otimes_{\mathrm{K}} \mathrm{E})\) -模）。
\item
  证明若 \([\mathrm{D}:\mathrm{K}]\) 有限且与 \(m\) 互素，则 \(\mathrm{D}\otimes_{\mathrm{K}}\mathrm{E}\) 是一个体（应用 \(a\)）。
\end{enumerate}

